\documentclass[11pt,reqno]{amsart}
\usepackage[a4paper,textwidth=6.2in,vmargin=0.92in]{geometry}
\usepackage{amsmath,amssymb}
\usepackage[T1]{fontenc}
\usepackage{lmodern}
\usepackage[hidelinks]{hyperref}
\allowdisplaybreaks[1]

\newtheorem{theorem}{Theorem}[section]
\newtheorem{proposition}[theorem]{Proposition}
\newtheorem{lemma}[theorem]{Lemma}
\theoremstyle{definition}
\newtheorem{definition}[theorem]{Definition}
\numberwithin{equation}{section}

\newcommand{\CC}{\mathbb C}
\newcommand{\QQ}{\mathbb Q}
\newcommand{\RR}{\mathbb R}
\newcommand{\ZZ}{\mathbb Z}
\newcommand{\re}{\operatorname{Re}}
\newcommand{\im}{\operatorname{Im}}
\newcommand{\res}{\operatorname{Res}}
\newcommand{\tr}{\operatorname{tr}}
\newcommand{\rd}{\operatorname{rd}}
\newcommand{\abs}[1]{\left|#1\right|}
\newcommand{\htle}[2]{\operatorname{ht}\bigl(#1\bigr)\le #2}
\newcommand{\HS}{\mathrm{HS}}
\newcommand{\of}{\mathcal O}

\title[Dirichlet zeros and theta-lift volumes]{The Dirichlet Riemann hypothesis via theta-lift covolumes}
\subjclass[2020]{Primary 11M26. Secondary 11F27, 11R42.}
\keywords{Dirichlet \(L\)-functions, theta lifts, covolumes, quadratic twists}
\hypersetup{pdftitle={The Dirichlet Riemann hypothesis via theta-lift covolumes}}

\begin{document}

\begin{abstract}
We present a proof of the Riemann hypothesis for Dirichlet
\(L\)-functions by comparing two estimates for the squared covolume
of a family of iterated theta lifts. Short integral forms and a
second moment give a family spanning almost all of the source.
For the lower bound, an auxiliary twist is chosen at each finite
place and the product formula is applied with a common untwisted
determinant. We use the arithmetic estimates of Chojecki for the
theta-lift construction of Alp\"oge. Cancellation of the common term
gives a one-sided estimate for a smoothed prime sum, from which
nonvanishing follows by quadratic twisting and Landau's argument.
\end{abstract}

\maketitle

\section{Introduction}

\begin{theorem}\label{thm:grh}
For every Dirichlet character \(\chi\), the function \(L(s,\chi)\)
is nonzero for \(\re s>1/2\), with the pole at \(s=1\) for a
principal character excepted. Every zero in \(0<\re s<1\)
has real part \(1/2\).
\end{theorem}

For a prime \(d\equiv7\bmod 8\), let \(\chi_d=(\,\cdot\,/d)\) and put
\(\Lambda_d(n)=\Lambda(n)(1-\chi_d(n))\), where \(\Lambda\) is the von
Mangoldt function. For \((q,d)=1\), define \(R_{q,d}\) by
\begin{equation}\label{eq:intro-remainder}
 \sum_{\substack{n\le x\\ n\equiv1\bmod q}}\Lambda_d(n)
       \left(\frac1n-\frac1x\right)
 = \frac{\log x}{\varphi(q)}+C_{q,d}+R_{q,d}(x).
\end{equation}
The constant \(C_{q,d}\) is characterized by \(R_{q,d}(x)\to0\).

\begin{theorem}\label{thm:prime}
Let \(q\) be a positive multiple of \(18480\), and let \(d\) be a
prime congruent to \(q/2-1\bmod q\). For every \(\epsilon>0\)
there is a constant \(A_{q,d,\epsilon}\) such that
\begin{equation}\label{eq:prime}
 R_{q,d}(x)\le A_{q,d,\epsilon} x^{-1/2+\epsilon}
 \qquad (x\ge2).
\end{equation}
\end{theorem}

The proof uses Alp\"oge's iterated theta-lift construction
\cite{alpoge} with the arithmetic theorems in Chojecki's
exposition \cite{chojecki}. The chain has length \(m\asymp\sqrt{x}\).
Local divisibility bounds the integral index of a family of short
forms. A second moment selects a small collection of translates
which span enough directions for a determinant comparison.
In the opposite direction, a product of critical Hecke values
controls the loss at each stage of the chain. The two bounds contain
the same contribution from the weights. Subtracting it leaves the
prime-sum remainder multiplied by \(m\).

We give the analytic implication in \S\ref{sec:analytic}, the data
and determinant comparison in \S\S\ref{sec:data}--\ref{sec:gram},
and the three estimates in \S\S\ref{sec:short}--\ref{sec:long}.
The last section combines them with the prime number theorem at the
larger of the two cutoffs.


\section{From a one-sided estimate to nonvanishing}\label{sec:analytic}

\subsection{The Mellin transform of the remainder}

Fix \(q\ge1\) and a prime \(d\equiv7\bmod 8\) with \((q,d)=1\).
Write \(A_{q,d}(u)=\sum_{n\le u,\ n\equiv1\bmod q}\Lambda_d(n)\) and
\(\mu_q=1/\varphi(q)\). Partial summation gives
\begin{equation}\label{eq:partial}
 \sum_{\substack{n\le x\\ n\equiv1\bmod q}}\Lambda_d(n)
       (n^{-1}-x^{-1})
 = \int_1^x A_{q,d}(u)\,\frac{du}{u^2}.
\end{equation}
Orthogonality of characters and the prime number theorem in arithmetic
progressions give, for every \(B>0\),
\[
 A_{q,d}(u)=\mu_q u+O_{q,d,B}\bigl(u(\log u)^{-B-1}\bigr)\qquad(u\ge2).
\]
Indeed the characters in the two sums are \(\psi\bmod q\) and
\(\psi\chi_d\bmod qd\). The second family has no principal member, since
the restriction to the \(d\)-factor is nontrivial. We use the prime number
theorem in the form of \cite[Ch.~11]{mv}. It follows that
\begin{equation}\label{eq:pnt-tail}
 C_{q,d}=\int_1^\infty\bigl(A_{q,d}(u)-\mu_q u\bigr)\frac{du}{u^2},
 \qquad R_{q,d}(x)\ll_{q,d,B}(\log x)^{-B}\quad\text{in absolute value}.
\end{equation}

For \(\re s>1\) define
\[
 F_{q,d}(s)
 =\frac{1}{\varphi(q)}\sum_{\psi\bmod q}
     \left(-\frac{L'}{L}(s,\psi)+\frac{L'}{L}(s,\psi\chi_d)\right).
\]
By the Euler products it is the Dirichlet series with coefficients
\(\mathbf 1_{n\equiv1\bmod q}\Lambda_d(n)\). For \(\re s>1\), integration
of the weight on \([n,\infty)\) gives
\(\int_n^\infty(n^{-1}-x^{-1})x^{-s}\,dx=n^{-s}/(s(s-1))\).
Consequently
\begin{equation}\label{eq:mellin}
 M_{q,d}(s):=\int_1^\infty R_{q,d}(x)x^{-s}\,dx
 = \frac{F_{q,d}(s)}{s(s-1)}
       -\frac{\mu_q}{(s-1)^2}-\frac{C_{q,d}}{s-1}.
\end{equation}
All the interchanges of sum and integral here are absolute. The pole
of \(F_{q,d}\) at \(1\) has residue \(\mu_q\). If \(f_0\) denotes its
constant Laurent coefficient there, the constant in
\eqref{eq:pnt-tail} is \(C_{q,d}=f_0-\mu_q\). This follows by multiplying
\eqref{eq:mellin} by \(s-1\) and letting \(s\downarrow1\), since
\eqref{eq:pnt-tail} makes \(M_{q,d}(s)\) bounded on approach from the
right. In particular,
\begin{equation}\label{eq:c1}
 C_{1,d}=-1-\gamma+\frac{L'}{L}(1,\chi_d),
\end{equation}
where \(\gamma\) is Euler's constant.

In the half-plane \(\re s>0\), the right side of \eqref{eq:mellin} is
meromorphic, regular at \(1\), and its only possible poles are simple.
Let \(m(\rho,\chi)\) denote the zero multiplicity of \(L(s,\chi)\) at
\(\rho\ne1\), with value \(0\) at a regular nonzero point. For
\(\re\rho>0\), \(\rho\ne1\), we have
\begin{equation}\label{eq:residue}
 \res_{s=\rho}M_{q,d}(s)
 =\frac{\displaystyle\sum_{\psi\bmod q}
                \bigl(m(\rho,\psi)-m(\rho,\psi\chi_d)\bigr)}
        {\varphi(q)\rho(1-\rho)}.
\end{equation}
The sign in this expression will be used at real \(\rho\).

\subsection{Selection of a quadratic twist}

\begin{lemma}\label{lem:twists}
Let \(8\mid q\) and let \(d_0\bmod q\) be a reduced residue class with
\(d_0\equiv 7\bmod 8\). Given \(1/2<\beta<1\) and \(H\ge0\), there is a
prime \(d\equiv d_0\bmod q\) such that all \(L(s,\psi\chi_d)\), with
\(\psi\bmod q\), are nonzero in \(\re s\ge\beta,\ |\im s|\le H\).
\end{lemma}

\begin{proof}
We give a smoothed mean-square argument. For \(j\in\{0,1\}\) let
\(w_j(v)=v^j e^{-\pi v^2}\). For \(|a_n|\le1\) set
\[
 \vartheta_{a,j}(u)=u^{1/2}\sum_{n\ge1}a_n w_j(nu),\qquad
 \gamma_j(s)=\tfrac12\pi^{-(s+j)/2}\Gamma((s+j)/2).
\]
Take \(X\ge2\), put \(T=(qX)^{1/2}\), and consider a nonprincipal
primitive character \(\chi\) of conductor \(Q\le T^2\) and parity \(j\).
Define
\[
 \begin{split}
 J(\chi)&=\int_{1/T}^\infty
    \bigl(|\vartheta_{\chi,j}(u)|+|\vartheta_{\mu\chi,j}(u)|\bigr)\frac{du}{u},\\
 B_\chi(s)&=\int_{1/T}^\infty
                \vartheta_{\mu\chi,j}(u)u^{s-1/2}\frac{du}{u},
 \end{split}
\]
where \(\mu\) is the M\"obius function. The theta transformation for
primitive characters gives
\(\abs{\vartheta_{\chi,j}(u)}
 =\abs{\vartheta_{\bar\chi,j}(1/(Qu))}\).
The inversion \(u\mapsto1/(Qu)\) and
\(\abs{\vartheta_{\chi,j}}=\abs{\vartheta_{\bar\chi,j}}\) therefore give
\[
 \int_0^\infty |\vartheta_{\chi,j}(u)|\,\frac{du}{u}
 =2\int_{Q^{-1/2}}^\infty|\vartheta_{\chi,j}(u)|\,\frac{du}{u}
 \le 2J(\chi).
\]
Mellin transformation gives
\(\int_0^\infty \vartheta_{\chi,j}(u)u^{s-1/2}du/u=\gamma_j(s)L(s,\chi)\).
The identity holds first by termwise integration for \(\re s>1\) and
extends by the theta transformation. The analogous integral for
\(\mu\chi\) equals \(\gamma_j(s)/L(s,\chi)\) for \(\re s>1\).
Since \(\vartheta_{a,j}(u)\ll u^{-1/2}\) for \(0<u\le1\), omitting
\((0,1/T)\) in the latter integral costs \(O(T^{-1})\) on \(\re s=2\).
Thus the function
\[
 G_\chi(s)=\gamma_j(s)L(s,\chi)B_\chi(s)-\gamma_j(s)^2
\]
is \(O(J(\chi)^2+1)\) on \(\re s=1/2\) and \(O(T^{-1})\) on
\(\re s=2\), with absolute implied constants. Its boundedness in the
closed strip, for fixed \(\chi,T\), follows from absolute integrability of
the Mellin kernels on both boundary lines.

Take \(0<\kappa<(\beta-1/2)/3\). If \(J(\chi)\le T^\kappa\), the
three-lines theorem bounds \(G_\chi(s)\) uniformly for
\(\beta\le\re s\le1\) by
\[
 O\bigl(T^{2\kappa(1-\theta)-\theta}\bigr),
 \qquad \theta=\frac{\re s-1/2}{3/2}.
\]
This tends to \(0\) as \(X\to\infty\). At a zero in
\(\beta\le\re s\le1,\ |\im s|\le H\), however,
\(\abs{G_\chi(s)}=\abs{\gamma_j(s)}^2\) is bounded below by a positive
constant depending only on \(\beta,H\). For all sufficiently large \(X\)
a character with a zero in this rectangle must satisfy \(J(\chi)>T^\kappa\).

We next bound \(J\) on average. Write \(n_{\rm odd}\) for the odd part
of \(n\) and, for \(v\ge1\), use the coefficient
\[
 a_n(v)=\alpha_n\left(\frac{-v}{n_{\rm odd}}\right),\qquad
 |\alpha_n|\le1,
\]
with the Jacobi symbol, including its value \(1\) at denominator \(1\).
For every \(\eta>0\) and \(1/T\le u\le1\),
\begin{equation}\label{eq:quad-second}
 \sum_{1\le v\le X}|\vartheta_{a(v),j}(u)|^2
   \ll_\eta X\log(2/u)+u^{-2-\eta}.
\end{equation}
To see this, expand the square. When \((nn')_{\rm odd}\) is a square,
bound the sum over \(v\) by \(X\). The number of these pairs with
\(n\le U,\ n'\le V\) is
\(\ll (UV)^{1/2}\log(2\min(U,V))\) for \(U,V\ge1\). Write
\(n=2^e h a^2,\ n'=2^{e'}h (a')^2\) with \(h\) odd and squarefree and
\(a,a'\) odd. Summing the bounds
\((UV)^{1/2}2^{-(e+e')/2}/h\) proves the count. Applying the count on
dyadic enlargements of \([1,1/u]^2\) and using the rapid decay of \(w_j\)
gives the first term of \eqref{eq:quad-second}.
For the remaining pairs, the character of \(v\) is nonprincipal.
P\'olya--Vinogradov, with inclusion--exclusion for an imprimitive modulus,
bounds its sum over \(v\) by \(O_\eta((nn')^{1/2+\eta/2})\).
Here a character of modulus \(M\) induced from conductor \(F\) costs at
most \(O(\tau(M)\sqrt{F}\log(2F))\), which suffices.
Since
\(\sum_n n^{1/2+\eta/2}|w_j(nu)|\ll_\eta u^{-3/2-\eta/2}\),
the factor \(u\) in the square of \(\vartheta\) gives the second term.
The argument uses only the usual character sum estimate
\cite[Ch.~9]{mv}. For \(u\ge1\), the corresponding sum of squares is
\(O(X e^{-c u^2})\) for an absolute \(c>0\).

Fix \(\psi\bmod q\) and let \(\psi^*\) be the primitive character inducing
it. For the primes \(d\) in the given class, \(\psi^*\chi_d\) is primitive
of conductor \(d\,\operatorname{cond}(\psi^*)\le qX\). The parity is fixed.
Quadratic reciprocity and \(d\equiv7\bmod8\) give
\((n/d)=(-d/n_{\rm odd})\). We apply \eqref{eq:quad-second} with
\(\alpha_n=\psi^*(n)\) and \(\alpha_n=\mu(n)\psi^*(n)\).
Cauchy--Schwarz, followed by integration in \(u\), gives
\[
 \sum_{\substack{d\le X\ {\rm prime}\\d\equiv d_0\bmod q}}
                     J(\psi^*\chi_d)\ll_{q,\eta} T^{2+\eta}
 \qquad(\eta>0).
\]
For example, apply \eqref{eq:quad-second} with exponent \(\eta\) on
\([1/T,1]\), where \(X\le T^2\) and \(u^{-2-\eta}\le T^{2+\eta}\),
and absorb \(\log T\) after Cauchy--Schwarz. The range \(u\ge1\) costs
\(O(X)\). With \(\eta=\kappa/2\), the number of excluded primes is
\(O_{q,\beta,H}(T^{2-\kappa/2})\), after summing over the finitely many
\(\psi\). The prime number theorem in the fixed progression provides
\(\asymp_q X/\log X\) primes, more than the excluded count. Finally the
Euler factors removed in passing from \(\psi^*\chi_d\) to
\(\psi\chi_d\) are nonzero in \(\re s>0\). This proves the lemma.
\end{proof}

\subsection{Positivity and the transform}

\begin{lemma}[Landau's argument]\label{lem:landau}
Let \(g\ge0\) be locally integrable on \([X,\infty)\), where \(X\ge1\),
and suppose \(\int_X^\infty g(x)x^{-s}\,dx\) converges for
\(\re s\) sufficiently large. Suppose it continues meromorphically
to \(\re s>\beta\) with at most simple poles and with nonpositive
residue at each possible real pole there. Then it is holomorphic in
\(\re s>\beta\).
\end{lemma}

\begin{proof}
Let \(a\) be the abscissa of convergence. If \(a>\beta\), the transform has
a singularity at \(a\). Indeed, if it were regular there, its Taylor
series at a real \(s_0>a\) sufficiently close to \(a\) would converge at
some \(s_0-h<a\). The series at that point has nonnegative terms
\(\int_X^\infty g(x)x^{-s_0}(h\log x)^j\,dx/j!\). Monotone convergence
would imply convergence of the integral at \(s_0-h\). The singularity
must therefore be a simple pole. Its nonzero residue must be positive,
since the integral is nonnegative as \(s\) approaches \(a\) from the
right. This contradicts the residue condition. Thus \(a\le\beta\),
and absolute convergence gives holomorphy in the required half-plane.
\end{proof}

\begin{proposition}\label{prop:criterion}
With \(q,d_0,\beta\) as in Lemma~\ref{lem:twists}, suppose that for each
prime \(d\equiv d_0\bmod q\) there is an eventual bound
\(R_{q,d}(x)\le A_{q,d,\beta}x^{\beta-1}\).
Then every \(L(s,\psi)\), \(\psi\bmod q\), is nonzero in
\(\re s>\beta\), with the pole at \(1\) excepted.
\end{proposition}

\begin{proof}
Take \(\rho\) with \(\beta<\re\rho<1\), and use Lemma~\ref{lem:twists}
with \(H=|\im\rho|\) to fix \(d\). For this \(d\) some \(B>0\) and \(X\)
make \(g(x)=B x^{\beta-1}-R_{q,d}(x)\) nonnegative for \(x\ge X\).
Its transform equals
\[
 \frac{B X^{\beta-s}}{s-\beta}-M_{q,d}(s)
       +\int_1^X R_{q,d}(x)x^{-s}\,dx .
\]
At a real \(s\in(\beta,1)\), the numerator in \eqref{eq:residue} is
a sum of untwisted multiplicities and the denominator is positive.
The residue of the transform of \(g\) is therefore nonpositive.
There are no possible poles on the real axis at or beyond \(1\).
Lemma~\ref{lem:landau} makes \(M_{q,d}\) regular at \(\rho\).
By \eqref{eq:residue} and the choice of \(d\), the sum of untwisted
multiplicities at \(\rho\) is \(0\), so every summand is \(0\).
The classical nonvanishing theorem on \(\re s=1\) and the Euler product
for \(\re s>1\) give the rest of the assertion.
\end{proof}

\begin{proof}[Proof of Theorem~\ref{thm:grh} from Theorem~\ref{thm:prime}]
For a character modulo \(q'\), use \(q=18480q'\) and \(d_0=q/2-1\).
Then \(d_0\) is coprime to \(q\) and \(d_0\equiv7\bmod 8\).
For every \(1/2<\beta<1\), Theorem~\ref{thm:prime} with
\(\epsilon=\beta-1/2\) gives the bound in
Proposition~\ref{prop:criterion}.
Inducing a character to modulus \(q\) only removes Euler factors
\(1-\chi(p)p^{-s}\), which are nonzero for \(\re s>0\).
Thus every \(L(s,\chi)\) is nonzero in \(\re s>1/2\), apart from
the stated pole. The primitive functional equation takes a zero in
the open strip at \(s\) to one at \(1-s\) for the dual character.
The gamma factors are regular and nonzero there. The finite Euler
factors give the claim for imprimitive characters as well.
\end{proof}




\section{The two cutoffs and the automorphic data}\label{sec:data}

Fix \(q,d\) as in Theorem~\ref{thm:prime} and \(0<\delta<1/100\).
All asymptotic statements below are for \(x\to\infty\) with \(q,d,\delta\)
fixed. We take
\[
 \log b=x^\delta,\qquad \log y=x^{\delta/6},\qquad
 \sqrt{x}\le m<2\sqrt{x},\qquad x^5\le r<3x^5,
\]
where \(m\) is a power of two and \(r\) a power of three. Write
\(k=r+1-m\), \(t=5r\), and set
\[
 s_0=22\log_2 m,\qquad c_0=\lfloor\log_3 s_0\rfloor+5,\qquad
 N_1=q\,2^{s_0}3^{c_0},\qquad
 E=\QQ(\sqrt{-d},\mu_{N_1}),\quad F=E^+.
\]
Because \(d\nmid N_1\), the quadratic field \(\QQ(\sqrt{-d})\) is
disjoint from \(\QQ(\mu_{N_1})\). Fixing the embedding of the quadratic
field identifies the complex places \(\Phi\) of \(E\) with
\((\ZZ/N_1\ZZ)^\times\). Write \(f=|\Phi|=[F:\QQ]\).
Let \(\mathcal T\) be the kernel of reduction to
\((\ZZ/q\ZZ)^\times\), of size \(T\), so \(f=\varphi(q)T\).
In the construction of \cite[Proposition~3.4]{chojecki} there is a set
\(\mathcal S\subset-\mathcal T\) of size \(S\) with \(T/S\le\log b\).
We take that choice of \(\mathcal S\), and take \(g=g_0g_1\) with
the two auxiliary split primes \(g_i\in[b^{10},b^{11}]\) from the
same proposition.

For \(z>0\) put \(\lceil z\rceil_{N_1}=N_1\lceil z/N_1\rceil\).
Use the order \(a=\lceil xb\rceil_{N_1}\), the nominal weights
\(\kappa_\tau=a/x\) for \(\tau\in\mathcal S\) and
\(\kappa_\tau=a/y\) otherwise, and the integral weights
\[
 w_j(\tau)=\lceil\kappa_\tau\rceil_{N_1}+(2j+1)r^8
 \qquad (0\le j\le m).
\]
Let \(w=f^{-1}\sum_{\tau\in\Phi}\log\kappa_\tau\) and
\(G=\log(g\,\rd(E)/(2\pi))\), where \(\rd(E)\) is the root
discriminant. The split level prime \(N\) of
\cite[Lemma~3.5]{chojecki} is chosen with \(\log N\) between
\(\lambda=m(w+G)-G-\tfrac{15}{2}\log b\) and
\(\lambda+\sqrt{\log b}\). In particular \(N\ge b^m\).
Fix the level generator \(\pi_N\) with \(\pi_N\overline{\pi_N}=N\)
used in that construction.

\subsection{The two-cutoff identity}

For \(\nu\) prime to \(N_1\), let \([-\nu]\in\Phi\) be the
residue class of \(-\nu\). If \((\nu,N_1)>1\) and
\(\Lambda(\nu)\ne0\), then \(\nu\) is a power of a prime dividing \(q\).
Every such prime splits in \(\QQ(\sqrt{-d})\), since
\(-d\equiv1\bmod p\) for odd \(p\mid q\) and \(-d\equiv1\bmod8\).
Thus \(\Lambda_d(\nu)=0\) in the nonunit case.
We can consequently assign any value to \(\kappa_{[-\nu]\tau}\)
there.

With \(v_+=\max(v,0)\), put
\[
 J_\tau=\sum_{\nu\le a}\Lambda_d(\nu)
             \left(\frac{a}{\nu}-\kappa_{[-\nu]\tau}\right)_+,
 \qquad
 \mathcal R=\sum_{\tau\in\mathcal T}J_\tau-
             Ta(\log a+C_{1,d}-w).
\]
For a unit \(\nu\), the number of \(\tau\in\mathcal T\) for which
\([-\nu]\tau\in\mathcal S\) is \(S\) when \(\nu\equiv1\bmod q\),
and zero otherwise. Indeed, multiplication by \([-\nu]\) maps
\(\mathcal T\) bijectively to the fibre over \(-\nu\bmod q\).
If \(P_{q,d}(z)\) denotes the prime sum on the left of
\eqref{eq:intro-remainder}, this gives
\[
 \sum_{\tau\in\mathcal T}J_\tau
 =aS\bigl(P_{q,d}(x)-P_{q,d}(y)\bigr)+aT P_{1,d}(y).
\]
Here \(a\ge y\), and taking the positive part imposes the relevant
cutoff. Substituting \eqref{eq:intro-remainder} and using
\(w=\log a-\frac{S}{f}\log x-(1-\frac{S}{f})\log y\) gives
\begin{equation}\label{eq:two-weights}
 \mathcal R=aS\bigl(R_{q,d}(x)-R_{q,d}(y)\bigr)+aT R_{1,d}(y).
\end{equation}

\subsection{The common leading term}

Set
\[
 \ell_d=\gamma-\frac{L'}{L}(1,\chi_d),\qquad
 \Omega=\frac{2\pi e^{\ell_d}}{\rd(E)},\qquad
 B=fa(\log b)^3,\qquad L=B\log b.
\]
The term shared by the two volume estimates is \(H=mku\), where
\begin{equation}\label{eq:direction}
 u=\sum_{\tau\in\Phi}\left[
  \log\bigl(w_0(\tau)!\bigr)-w_0(\tau)\log(\Omega/g)
  -\sum_{\nu\le a}\Lambda_d(\nu)
          \left(\frac{\kappa_\tau}{\nu}-\kappa_{[-\nu]\tau}\right)_+
 \right].
\end{equation}
This is the direction cost with the normalization
of \cite[{\rm(3.8)}]{chojecki}.
We will only need the crude absolute bound
\begin{equation}\label{eq:coarse-main}
 |H|+m|\mathcal R|\le b^2x^{99}.
\end{equation}
Indeed \(N_1=O_q(x^{11}\log x)\), \(xb\le a\le 2xb\) for
large \(x\), and \(w_0(\tau),\kappa_\tau=O(b)\) uniformly in
\(\tau\). Also \(\log\rd(E)=O_{q,d}(\log x)\) by the cyclotomic
discriminant bound and the inclusion \(E\subset\QQ(\mu_{dN_1})\).
The factorial and period terms in \eqref{eq:direction} are
\(O_{q,d}(b\log a)\) per place, and the prime sum is bounded by
\[
 2\kappa_\tau\sum_{\nu\le a}\frac{\log\nu}{\nu}
       \ll b(1+\log a)^2 .
\]
As \(f\le N_1\), these estimates give \(|H|=O_{q,d,\delta}(bx^{21})\).
The remainders in \eqref{eq:two-weights} are bounded by
\eqref{eq:pnt-tail}, so \(m|\mathcal R|=O_{q,d,\delta}(maT)\).
Both bounds imply \eqref{eq:coarse-main} for large \(x\).

\subsection{Spaces, translates, and lifts}

We use the unitary spaces, finite levels, and theta data of
\cite[\S\S4--6 and Theorem~20.16]{chojecki} over \(E/F\).
Let \(X_j\) be the definite spaces of ranks \(k+j(t+1)\), and
\(Y_j\) the Hermitian spaces of ranks \(r+jt\) in the chain,
with \(Y_m\) definite. Write \(\mathcal V\) for the space of
algebraic modular forms on \(U(X_0)\) with the indicated Schur
coefficients. At the marked places over \(g_0,g_1\) its level is
principal, and at each \(w\mid\pi_N\) it fixes the last row modulo
\(w\). Thus a form is a right-level-invariant function on
\(U(X_0)(\mathbb A_{F,f})\), transforming under \(U(X_0)(F)\) by
the coefficient representation. Its values on the finite double
quotient satisfy the corresponding stabilizer conditions.

Let \(\Gamma\) be the product of the groups
\(\operatorname{GL}_r(k_w)\) at the marked split places over
\(g_0,g_1\), and let \(\rho\) be its unitary action by finite
translation on the principal-level cusp forms \(\mathcal P\)
on \(U(Y_0)\) of the prescribed discrete-series type. Write \(e_0\)
for the central projector in the construction and
\(\mathcal W=e_0\mathcal P\). The projector retains the finite
types with the specified squarefree Green labels~\cite{green} and a
cuspidal label over each of \(g_0,g_1\).
The dimensions of the full column space and the projected row
space are denoted by \(C=\dim\mathcal V\) and \(D=\dim\mathcal W\),
respectively. In particular \(C>0\) by the free-level class
description, since the coefficient representation is nonzero.

At \(w\mid\pi_N\), let \(z_w\) denote translation by a scalar
lifting a generator of \(k_w^\times\). Put
\begin{equation}\label{eq:ramified}
 E_N=\prod_{w\mid\pi_N}(1-z_w),\qquad
 \widehat{\mathcal V}=E_N\mathcal V.
\end{equation}
Let \(P:\mathcal V\to\mathcal P\) be the Poincar\'e series map
formed using the Bessel datum at infinity and the prescribed
finite column data, in the normalization of
\cite[\S11]{chojecki}. The row maps before rational rescaling
are
\begin{equation}\label{eq:row-maps}
 A_s=c_*e_0\rho(s)P\big|_{\widehat{\mathcal V}},
 \qquad s\in\Gamma,
\end{equation}
where \(c_*\ne0\) includes the common seesaw scalar and the
factor \(|\Gamma|\). These maps use the same datum for the
short-form and second-moment theorems, as specified in
\cite[Theorem~20.16]{chojecki}.

The twists \(\xi\) are the weight-zero tame characters in this
construction, whose first \(m\) powers are nontrivial at every
place of their support. The support is disjoint from its conjugate,
and its conductor norm \(q_\xi\) is prime to \(gN\).
The trivial twist has \(q_1=1\). Use the composite of the \(m\)
theta lifts from \(U(Y_0)\) to \(U(Y_m)\), with the prescribed
splittings and powers of \(\xi\), to obtain an injection
\[
 \Theta_\xi:\mathcal W\longrightarrow\mathcal H_\xi .
\]
Here \(\mathcal H_\xi\) is the terminal space of definite forms
with its squared class-sum norm. The injectivity follows from
the triangle identities and the inclusion-datum lifts in
\cite[\S13.8]{chojecki}.
At a finite embedding, the integral forms have values in the
transported integral Weyl lattices at each finite adelic point.
The field, the spaces, the levels, the integral structures, and
the lifts from this construction are fixed throughout
\S\S\ref{sec:gram}--\ref{sec:long}.




\section{Integral indices and determinant ratios}\label{sec:gram}

For each rational prime \(p\) we use the field isomorphisms
\(\iota:\CC\to\overline{\QQ}_p\), with
\(\abs{p}_\iota=p^{-1}\). Write \(\of_\iota\) for the inverse
image of the valuation ring and \(\bar\iota=\iota\circ c\), where
\(c\) denotes complex conjugation. These isomorphisms are
abstract field maps. Every embedding of a number field into
\(\overline{\QQ}_p\) extends to such a map.

For a twist with \(p\nmid q_\xi\), let
\(\mathcal I_{\xi,\iota}\subset\mathcal H_\xi\) be its integral
module. We use the class-sum pairing with conjugate gauges,
\begin{equation}\label{eq:int-pairing}
 \abs{\langle v,w\rangle}_\iota\le1
 \qquad
 \bigl(v\in\mathcal I_{\xi,\bar\iota},\
       w\in\mathcal I_{\xi,\iota}\bigr).
\end{equation}
Indeed, the endpoint has free principal level. In class
representatives integral above \(p\), the pairing uses the
integral Weyl Gram matrices and is a sum without stabilizer
denominators. These are the integral realizations of
\cite[\S6.13, Lemma~8.12, and {\rm(14.3)}]{chojecki}.
At \(\xi=1\) the fixed Weyl coordinates on the classes give a
positive definite rational Gram matrix. Write \(v^\sigma\) for
the vector obtained by applying \(\sigma\in\operatorname{Aut}(\CC)\)
to these coordinates, and \(\|v\|_\sigma=\|v^\sigma\|\).

\begin{definition}
A family \(\mathbf v=(v_1,\ldots,v_n)\) of rows has index at
most \(I\) if there are \(a_p\ge0\) with \(\sum_p a_p\le I\)
such that, for every \(\iota\mid p\) and \(p\nmid q_\xi\),
some \(U\in M_n(\of_\iota)\) satisfies
\[
 |\det U|_\iota\ge e^{-a_p},\qquad
 \Theta_\xi\left(\sum_{j=1}^n U_{ij}v_j\right)
       \in\mathcal I_{\xi,\iota}\quad (1\le i\le n).
\]
The numbers \(a_p\) are uniform in \(\xi,\iota\).
The matrix may depend on both. The squared covolume at
\(\sigma\) is
\[
 \Delta_\sigma(\mathbf v)=
 \det\left(\left\langle
       (\Theta_1 v_i)^\sigma,(\Theta_1 v_j)^\sigma
     \right\rangle\right)_{i,j=1}^n .
\]
\end{definition}

\begin{definition}
For a family \(z=(z_\iota)_\iota\), the bound
\(\htle{z}{A}\) means that there are \(d_0\in\CC^\times\),
\(\alpha\in\RR\), and real \(a_p\) with
\[
 |\sigma d_0|\le e^\alpha,\qquad
 |d_0z_\iota|_\iota\le e^{a_p}\quad(\iota\mid p),\qquad
 \alpha+\sum_{p<Y}a_p\le A
 \quad\text{for all sufficiently large \(Y\)}.
\]
\end{definition}

Bounds add under multiplication, with the product of the
two witnesses as witness. A local unit at each \(\iota\)
can be inserted in the family without changing a bound.
If \(\htle{(1)_\iota}{A}\), then \(A\ge0\).
Indeed, a nonzero complex number with bounded conjugates
under all automorphisms of \(\CC\) is algebraic. Otherwise,
an isomorphism between rational function fields taking it
to an arbitrarily large transcendental number extends to an
automorphism of \(\CC\). Apply the rational product formula
to the norm of \(d_0\) from \(\QQ(d_0)\). The log of the
ordinary or \(p\)-adic absolute value of the norm, divided
by \([\QQ(d_0):\QQ]\), is bounded by \(\alpha\) or \(a_p\),
respectively. For \(Y\) containing all the primes in that norm,
the product formula gives the assertion.

Choose any Hermitian metric on \(\mathcal W\), the same at
every twist, and put \(Z_\xi=\det(\Theta_\xi^*\Theta_\xi)>0\).
Ratios of the \(Z_\xi\) are independent of the metric.

\begin{lemma}\label{lem:gram}
Let \(\mathbf v\) be an independent family of \(n\) rows
of index at most \(I\), with \(\Delta_\sigma(\mathbf v)\le e^{V_0}\)
for all \(\sigma\). Complete it to a basis by \(j=D-n\)
vectors \(w_a\) with \(\|\Theta_1 w_a\|_\sigma^2\le e^{A_0}\).
Choose one twist \(\xi_\iota\) for each \(\iota\mid p\) with
\(p\nmid q_{\xi_\iota}\).
Suppose there are nonzero \(\gamma_\iota\) such that
\(\gamma_\iota\Theta_{\xi_\iota}w_a\) belongs to both
\(\mathcal I_{\xi_\iota,\iota}\) and
\(\mathcal I_{\xi_\iota,\bar\iota}\) for every \(a\), and
\(\htle{(|\gamma_\iota|^{-2})_\iota}{G_0}\). Then
\begin{equation}\label{eq:gram-family}
 \htle{(Z_{\xi_\iota}/Z_1)_\iota}{V_0+2I+j(A_0+G_0)}.
\end{equation}
\end{lemma}

\begin{proof}
Let \(D_\xi\) be the Gram determinant of the images of the
full basis at \(\xi\). The change of basis gives
\(D_\xi/D_1=Z_\xi/Z_1\). To bound the finite values at
\(\iota\), use the index matrices at \(\iota\) and
\(\bar\iota\), both for the twist \(\xi_\iota\).
On the completion vectors use the scalar \(\gamma_\iota\).
By \eqref{eq:int-pairing} the new pairing matrix is integral
at \(\iota\). Taking its determinant and dividing by the
two index determinants yields
\[
 \bigl|(\overline{\gamma_\iota}\gamma_\iota)^j
             D_{\xi_\iota}\bigr|_\iota\le e^{2a_p}.
\]

The rational Gram matrix in the fixed coordinates satisfies
\(\sigma\langle v,w\rangle=\langle v^{c\sigma c},w^\sigma\rangle\).
It follows that \(\sigma D_1\) is the Hermitian pairing of
two exterior products, formed using \(c\sigma c\) and
\(\sigma\), respectively. For either exterior product the
squared norm is at most \(e^{V_0+jA_0}\). This follows by
adjoining the completion vectors one at a time and bounding
the norm of each orthogonal projection by the full norm.
Cauchy--Schwarz now gives \(|\sigma D_1|\le e^{V_0+jA_0}\).
Thus \(D_1\) witnesses the bound \(V_0+jA_0+2I\) for the
family \(\bigl((\overline{\gamma_\iota}\gamma_\iota)^j
Z_{\xi_\iota}/Z_1\bigr)_\iota\).
Multiplying by the inverse scalar family gives
\eqref{eq:gram-family}.
\end{proof}




\section{Short forms and their index}\label{sec:short}

Let \(\mathcal M_\iota\subset\mathcal V\) consist of columns
integral in the transported Weyl lattices at \(\iota\), and put
\(\mathcal M=\widehat{\mathcal V}\cap\bigcap_\iota\mathcal M_\iota\).
The quantitative local estimates of
\cite[Proposition~9.14, {\rm(5.1)--(5.3)}]{chojecki} give
a set \(\mathcal S_0\subset\mathcal M\) spanning
\(\widehat{\mathcal V}\), finitely supported nonnegative
\(g_p,e_p\), and \(h_0\) with
\begin{equation}\label{eq:short-estimates}
 \begin{gathered}
 \sum_p e_p\log p\le B,\qquad h_0\le H-m\mathcal R+B,\\
 \|\Theta_1 A_s\phi\|_\sigma^2
   \le \exp\left(h_0+2\sum_p g_p\log p+2B+x^{99}\right)
       \qquad (s\in\Gamma,\ \phi\in\mathcal S_0).
 \end{gathered}
\end{equation}
The local part of the estimates has the following form. For
each \(p,\iota\mid p\), and integer \(\beta<g_p-e_p+1\),
a finite family in \(\mathcal M_\iota\) can be recombined
over \(\of_\iota\) with determinant of absolute value at least
\(p^{-e_p C}\) into a submodule \(\mathcal N\) satisfying
\begin{equation}\label{eq:short-gain}
  p^{-\beta}\Theta_\xi A_s(\mathcal N\cap\widehat{\mathcal V})
       \subset\mathcal I_{\xi,\iota}
       \qquad (s\in\Gamma,\ p\nmid q_\xi).
\end{equation}
These estimates use the common seesaw normalization in
\eqref{eq:row-maps}. The local loss comes from the sublattice
on which the two halves of the Hecke repair can be combined,
and from the repair at the level prime \(N\). The factor \(C\)
in the determinant bound is the full column dimension.

\begin{proposition}\label{prop:short}
There is a common positive rational \(q_0\), with
\(B_s=q_0A_s\), such that
\[
 \|\Theta_1 B_s\phi\|_\sigma^2\le e^{H-m\mathcal R+L}
 \qquad(s\in\Gamma,\ \phi\in\mathcal S_0).
\]
For any subset \(\mathcal U\subset\Gamma\), a finite family
in \(\bigcup_{s\in\mathcal U} B_s(\mathcal M)\) has index at
most \(|\mathcal U|CB\).
\end{proposition}

\begin{proof}
Put \(\beta_p=\lceil g_p-e_p\rceil\) and take
\(q_0=\prod_p p^{-\beta_p}\).
The bound on \(h_0\) and
\(\log q_0\le\sum_p(e_p-g_p)\log p\) give
\[
 \|\Theta_1 B_s\phi\|_\sigma^2
   \le \exp\bigl(H-m\mathcal R+5B+x^{99}\bigr)
   \le e^{H-m\mathcal R+L}.
\]
The last inequality uses \(5B+x^{99}\le B\log b\) for large
\(x\), since \(a\ge xb\).

For each row in the given family choose one preimage in
\(\mathcal M\), and group the indices according to \(s\).
At \(\iota\mid p\), apply the local recombination with
\(\beta_p<g_p-e_p+1\) to the preimages in each group. The
recombined columns are still in \(\widehat{\mathcal V}\).
The factor \(q_0p^{\beta_p}\) is a unit of \(\of_\iota\),
so \eqref{eq:short-gain} gives integral lifts for the
recombined rows. The block diagonal recombination has
determinant of absolute value at least \(p^{-|\mathcal U|Ce_p}\).
The sum of the resulting logarithmic index allowances is
at most \(|\mathcal U|CB\) by \eqref{eq:short-estimates}.
\end{proof}




\section{The second moment and spanning}\label{sec:span}

For a subspace \(J\subset\mathcal V\), let \(\Pi_J\) denote its
orthogonal projector in the column norm used for \(P\).
The moment and projection estimates of
\cite[{\rm(11.13), (11.14), (11.18)}]{chojecki} give a common
normalizing factor \(M_*>0\) with
\begin{align}
 D&\le |\Gamma|M_* ,\label{eq:mass-dim}\\
 \sum_{s\in\Gamma}\|P^*\rho(s)P\|_{\HS}^2
      &\le (1+9b^{-4})\, C^2/M_* ,\label{eq:sharp-moment}\\
 \|e_0P\Pi_J\|_{\HS}^2
      &\ge \dim J-5b^{-4}C .\label{eq:retained-proj}
\end{align}
Here \(\|\cdot\|_{\HS}\) is the Hilbert--Schmidt norm.
The factor \(M_*\) is the product of the real formal degrees
and the mass of a definite genus of rank \(r\).
The flag integral estimates and the removal of the open-cell
restriction in \cite[Propositions~10.11 and~11.11]{chojecki}
give \eqref{eq:sharp-moment} for the full sum. Its leading
constant agrees with the one in the dimension bound
\eqref{eq:mass-dim}. Taking \(J=\mathcal V\) in
\eqref{eq:retained-proj} gives \(D>0\).

We also have
\begin{equation}\label{eq:ram-kernel}
 \dim\ker E_N\le \frac{f}{N-1}C.
\end{equation}
For this, a nonidentity power of \(z_w\) moves each level class.
In fact, a fixed class would yield a lattice automorphism
congruent to the identity at the marked principal level, which
is trivial by \cite[Lemma~8.12]{chojecki}. The level at \(w\)
fixes the last row, so a scalar belongs to it only when its
reduction is \(1\). Since \(N\) splits completely, the action
of \(z_w\) has cycles of length \(N-1\). A \(z_w\)-invariant
form is determined by one value on each cycle, giving
\(\dim\ker(1-z_w)\le C/(N-1)\). The rank inequality for a
product of operators and the \(f\) places over \(\pi_N\) give
\eqref{eq:ram-kernel}.

\begin{proposition}\label{prop:span}
There is a subset \(\mathcal U\subset\Gamma\) such that
\[
 |\mathcal U|C\le5D\log(1+b^2),\qquad
 \dim\sum_{s\in\mathcal U}B_s(\widehat{\mathcal V})
       \ge(1-4b^{-2})D.
\]
\end{proposition}

\begin{proof}
For \(s\in\Gamma\) set \(Q_s=e_0\rho(s)P\) on
\(\mathcal V\). Given a subspace \(F_0\subset\mathcal W\)
of dimension \(\alpha D\), consider the positive operator
\[
 X_*=\frac1{|\Gamma|}\sum_{s\in\Gamma}
               \rho(s)PP^*\rho(s)^*
 \qquad \text{on }\mathcal P .
\]
Taking the trace of \(X_*^2\) and using the group action gives
\[
 \|X_*\|_{\HS}^2
   =\frac1{|\Gamma|}\sum_{s\in\Gamma}\|P^*\rho(s)P\|_{\HS}^2
   \le\frac{(1+9b^{-4})C^2}{D}.
\]
For some \(s\) the trace
\(\|\Pi_{F_0}\rho(s)P\|_{\HS}^2\) is at most its average
\(\tr(\Pi_{F_0}X_*)\). Cauchy--Schwarz bounds this by
\[
 C\sqrt{\alpha(1+9b^{-4})}
       \le C\left(\frac{1+\alpha}{2}+\frac92b^{-4}\right).
\]
On \(J=Q_s^{-1}(F_0)\), the operator \(Q_s\) agrees with
\(\Pi_{F_0}\rho(s)P\). By \eqref{eq:retained-proj}
and commutation of \(e_0\) with \(\rho(s)\),
\[
 \dim J-5b^{-4}C
 \le\|e_0P\Pi_J\|_{\HS}^2
 \le\|\Pi_{F_0}\rho(s)P\|_{\HS}^2.
\]
For the inverse image under \(Q_sE_N\), the dimension is
at most \(\dim J+\dim\ker E_N\). Therefore for some \(s\)
\[
 \dim (Q_sE_N)^{-1}(F_0)
       \le C\left(\frac{1+\alpha}{2}+\eta\right),
 \qquad
 \eta=b^{-2}.
\]
Here \(\frac{19}{2}b^{-4}+f/(N-1)\le\eta\) for large \(x\),
since \(f=O_q(x^{11}\log x)\) and \(N\ge b^m\).
The image of \(Q_sE_N\) is exactly \(B_s(\widehat{\mathcal V})\)
by \eqref{eq:ramified} and \eqref{eq:row-maps}.

Start with \(F_0=0\). Whenever its codimension \(z\) in
\(\mathcal W\) exceeds \(4\eta D\), choose \(s\) as above
and add the image of \(Q_sE_N\).
The increase in dimension is
\(\Delta\ge C(z/(2D)-\eta)>0\), by the rank of the map to
\(\mathcal W/F_0\). In particular no index is chosen twice.
The decrease of \(5D\log(1+z/(\eta D))\) in this step is at least
\[
 \frac{5D\Delta}{z+\eta D}\ \ge\
 \frac{5Cz}{4(z+\eta D)}\ \ge\ C.
\]
The process stops with codimension at most \(4\eta D\), and
the initial value of the logarithm gives the asserted bound
on \(|\mathcal U|C\).
\end{proof}




\section{A lower bound for the lifted volume}\label{sec:long}

The determinants at a single stage of the chain can be
computed on a source space of fixed dimension. The necessary
relation with the composite lift uses the measures and
integral gauges from \S\ref{sec:data}. We give the determinant
argument after stating the comparison and specialization
results in these normalizations.

\subsection{Comparison at one stage}

For each \(j\), \(1\le j\le m\), take the single lift
\(u_{j,\xi}\) in the definite comparison
of \cite[\S13.9]{chojecki}. Its source \(U_j\) has dimension
\(D\), and \(\xi\) in this notation includes the prescribed
\(j\)-th power. Write
\[
 L_{j,\xi}=\operatorname{im}u_{j,\xi},\qquad
 \lambda_{j,\xi}=\det(u_{j,\xi}^*u_{j,\xi}).
\]
These lifts are injective. By
\cite[Proposition~14.10]{chojecki} there are maps
\(g_{j,\xi}\) from a fixed Hermitian space of tables
\(\mathcal T_j\) to the target of \(u_{j,\xi}\), and a fixed
surjection \(A_j:\mathcal T_j\to U_j\), with
\begin{equation}\label{eq:comparison-adj}
 u_{j,\xi}^*g_{j,\xi}=\gamma_{j,\xi} A_j,\qquad
 \left|\frac{\gamma_{j,\xi}}{\gamma_{j,1}}\right|^2
          =\frac{\Gamma_j(\xi)}{\Gamma_j(1)}.
\end{equation}
The nonzero scalars \(\gamma_{j,\xi}\) result from three
complementary-rank theta exchanges. The tables encode
the values in the small definite block. The separation
of the good forms by the product periods
\cite[Proposition~14.21]{chojecki} gives the surjection.

Let \(h_j>0\) be the fixed Weyl hook factor of the small
block and compress the comparison operator to the lift:
\[
 C_{j,\xi}=h_j\,\operatorname{pr}_{L_{j,\xi}}\,
                   g_{j,\xi}g_{j,\xi}^*\big|_{L_{j,\xi}}.
\]
Since \(g_{j,\xi}^*u_{j,\xi}=\overline{\gamma_{j,\xi}}A_j^*\)
is injective, \(C_{j,\xi}\) is positive definite. On \(U_j\),
\(u_{j,\xi}^*C_{j,\xi}u_{j,\xi}=h_j|\gamma_{j,\xi}|^2 A_jA_j^*\).
Taking determinants in orthonormal bases gives
\[
 \lambda_{j,\xi}\det C_{j,\xi}
   = h_j^D|\gamma_{j,\xi}|^{2D}\det(A_jA_j^*).
\]
Consequently
\begin{equation}\label{eq:comparison-ratio}
 \frac{\lambda_{j,1}}{\lambda_{j,\xi}}
 = \frac{\det C_{j,\xi}}{\det C_{j,1}}
        \left(\frac{\Gamma_j(1)}{\Gamma_j(\xi)}\right)^D.
\end{equation}

\subsection{A choice of twists for all stages}

Put \(B_0=B/(mk)\).
We use the choice of \(\xi_\iota\), for every \(\iota\mid p\),
from \cite[Proposition~20.15]{chojecki}. It has \(p\nmid q_{\xi_\iota}\).
The simultaneous specialization preserves the trace denominators
for the single lifts and the rank-one tests, and the valuations
of the required ramified critical values, for all powers through
\(m\). In particular the same \(\xi_\iota\) works in
\eqref{eq:comparison-ratio} at every stage.

Let \(e\) denote the finite projector on the target of the
single lift, and \(T_w\) the spherical Hecke polynomial
evaluated at the first root of its lift string at a split
place \(w\) off the levels. The image \(L_{j,\xi}\) is in
\(\operatorname{im}e\cap\ker T_w\). The unwanted space at \(w\)
is \(\operatorname{im}e\cap\ker T_w\cap L_{j,\xi}^\perp\).
We call a form tested if it is orthogonal to this space
for all such \(w\). The period vanishing in
\cite[Propositions~14.10 and~19.17]{chojecki} says that
the forms in \(g_{j,\xi}A_j^*U_j\) are tested.

Fix \(j,\iota\mid p\) and \(\xi=\xi_\iota\), and abbreviate
\(g=g_{j,\xi}\), \(h=h_j\), \(L'=L_{j,\xi}\).
The integral gauges in \cite[\S14.5]{chojecki} give tables \(f_a\),
integral at every embedding above \(p\), with
\(\langle f_a,\zeta\rangle=h q_a(\zeta)\), where \(q_a\) are the
integral coordinates of the table. The factor \(h\) clears
the inverse of its Weyl Gram matrix. The linear functionals
\(\ell_a(z)=\langle g f_a,z\rangle\) span \((L')^\vee\)
by injectivity of \(g^*|_{L'}\).
From their evaluations on any basis choose a \(D\)-row minor
of maximal \(\iota\)-adic absolute value. In the basis \(z_b\)
dual to the selected functionals, Cramer's rule gives
\(|\ell_a(z_b)|_\iota\le1\) for every \(a,b\).
As \(q_a(hg^*z_b)=\ell_a(z_b)\), the table \(hg^*z_b\) is
integral at \(\iota\). The corresponding form \(g(hg^*z_b)\)
is integral there and is tested by \eqref{eq:comparison-adj}.
In the basis \(z_b\) the entries of \(C_{j,\xi}\) are
\[
 \left\langle g f_a,\ \operatorname{pr}_{L'}\,g(hg^*z_b)\right\rangle,
\]
with \(a\) ranging over the selected indices. The tested
projection bound of \cite[Proposition~19.13]{chojecki},
applied with the trace relations of Proposition~19.7 there,
gives
\begin{equation}\label{eq:compressed-finite}
 |\det C_{j,\xi_\iota}|_\iota\le e^{Dr\beta_{j,p}},
 \qquad \sum_p\beta_{j,p}\le B_0,\quad\beta_{j,p}\ge0 .
\end{equation}

At \(\xi=1\) use the rational coordinates for the class-sum
and table metrics. The comparison norm bound in
\cite[Proposition~14.10]{chojecki} gives
\(h_j\|g_{j,1}v\|_\sigma^2\le e^{2rB_0}\|v^\sigma\|^2\).
For tested vectors the projection to \(L_{j,1}\) contracts
each conjugated norm. Indeed, at a fixed \(w\) the orthogonal
complement of the unwanted space is
\(L_{j,1}+\ker e+\operatorname{im}T_w\).
The last two summands remain perpendicular to \(L_{j,1}\)
in the conjugated gauges, by self-adjointness of \(e\)
and normality of the spherical test there
\cite[{\rm(14.38)}]{chojecki}.
If \(C_{j,1}z=\alpha z\) with \(z\ne0\), write
\(\zeta=g_{j,1}^*z\ne0\) and choose a table \(v\) with
\(v^{c\sigma c}=\zeta^\sigma\).
Then
\[
 \big|\sigma\langle g_{j,1}v,C_{j,1}z\rangle\big|
 =|\sigma\alpha|\,\|\zeta^\sigma\|^2
 \le h_j\|g_{j,1}v\|_{c\sigma c}\|g_{j,1}\zeta\|_\sigma
 \le e^{2rB_0}\|\zeta^\sigma\|^2 .
\]
The first inequality uses Cauchy--Schwarz and the contraction
on the tested form \(g_{j,1}\zeta\).
Taking the product over eigenvalues gives
\(|\sigma\det C_{j,1}|\le e^{2DrB_0}\).
Together with \eqref{eq:compressed-finite}, this shows that
\(\det C_{j,1}\) witnesses
\[
 \htle{\bigl(\det C_{j,\xi_\iota}/\det C_{j,1}\bigr)_\iota}{3DrB_0}.
\]


For the other factor of \eqref{eq:comparison-ratio}, let
\(z_0=(k+7m)/2\), \(z_+=m\), \(z_-=m/2\).
The exact product in \cite[{\rm(14.39)--(14.41)}]{chojecki}
has one block of critical values of each of these lengths.
In terms of the positive critical-value ratios \(\rho_{j,i,a}\)
of that product, it gives
\begin{equation}\label{eq:hecke-ratio}
 \frac{\Gamma_j(1)}{\Gamma_j(\xi)}
 =q_\xi^{-(z_0^2+z_+^2+z_-^2)}
    \prod_{i\in\{0,+,-\}}\prod_{a=1}^{z_i}\rho_{j,i,a}(\xi)^{-2}.
\end{equation}
Here a ratio is the absolute value of the twisted imprimitive
Hecke value divided by its untwisted counterpart, at \(a\)
with the character fixed in its block. The block \(-\) uses
the inverse twist. By
\cite[Propositions~14.23 and~20.15]{chojecki},
\(\bigl(\rho_{j,0,a}(\xi_\iota)^{-1}\bigr)_\iota\) has height
at most \(6B_0-u\) for \(a>7m/2\), and every other inverse
ratio has height at most \(x^{999}\).
The first estimate comes from the rank-one period tests.
For the other factors, the ramified-value interpolation
includes \(a=1\). The conductor term in
\eqref{eq:hecke-ratio} is a local unit at the chosen embeddings.
There are \(k/2\) factors with the first height bound and
\(5m\) with the second, before squaring. Hence
\[
 \htle{\bigl(\Gamma_j(1)/\Gamma_j(\xi_\iota)\bigr)_\iota}
       {k(6B_0-u)+10mx^{999}}.
\]
Combining this with \eqref{eq:comparison-ratio} and
the height bound for the compressed determinant gives
\begin{equation}\label{eq:single-lower}
 \htle{\bigl(\lambda_{j,1}/\lambda_{j,\xi_\iota}\bigr)_\iota}
       {Dk\bigl((\log b)B_0-u\bigr)}.
\end{equation}
We used \(10mx^{999}\le kB_0\) and \(7k+3r\le k\log b\),
both of which hold for large \(x\).

\subsection{The determinant of the chain}

The chain identity \cite[{\rm(13.55)}]{chojecki} gives positive
\(\kappa_\xi\) with
\[
 Z_\xi=\kappa_\xi^D\prod_{j=1}^m\lambda_{j,\xi},
 \qquad
 \htle{\bigl((\kappa_1/\kappa_{\xi_\iota})^2\bigr)_\iota}
       {2m^2t x^{999}}\le 2L .
\]
The equality is for the Gram determinant of the chain with
the class-sum pairing at its endpoint. Keeping the same scalar
as in this pairing is necessary in comparing the squared volumes.
Adding \eqref{eq:single-lower} for all the stages gives
\(D(L-H)\) for the bound on the height of their product of
reciprocal determinant ratios. The chain scalar therefore gives
\begin{equation}\label{eq:chain-lower}
 \htle{\bigl((Z_1/Z_{\xi_\iota})^2\bigr)_\iota}{2D(2L-H)} .
\end{equation}

The completion result \cite[Proposition~13.34]{chojecki}
provides a spanning family of \(\mathcal W\) satisfying the
completion conditions in Lemma~\ref{lem:gram} with
\(A_0=2y^2b^2\) and \(G_0=4m^2t x^{999}\).
It uses one integral scalar for each twist at all embeddings
prime to the conductor. In passing to \(\mathcal W\), the
operator \(|\Gamma|e_0\) acts by integral finite translations
on the chain. The extra factor in its norm is included in
\(A_0\). We have \(A_0+G_0\le b^3/4\) for large \(x\).

\begin{proposition}\label{prop:long}
If an independent family of \(n\ge(1-4b^{-2})D\) rows has
index at most \(I\) and squared covolume at most \(e^{V_0}\)
at every \(\sigma\), then
\begin{equation}\label{eq:lower}
 D(H-3L)\le V_0+2I.
\end{equation}
\end{proposition}

\begin{proof}
Complete the family by \(D-n\) elements of the spanning
family above. Lemma~\ref{lem:gram} gives
\[
 \htle{(Z_{\xi_\iota}/Z_1)_\iota}{V_0+2I+Db}
\]
because \((D-n)(A_0+G_0)\le Db\).
The square of this family times the family in
\eqref{eq:chain-lower} is \(1\). The product formula in
\S\ref{sec:gram} gives
\(0\le 2(V_0+2I+Db)+2D(2L-H)\).
Since \(b\le L\), this proves \eqref{eq:lower}.
\end{proof}




\section{The prime-sum bound}\label{sec:end}

\begin{proof}[Proof of Theorem~\ref{thm:prime}]
Choose \(\mathcal U\) by Proposition~\ref{prop:span}.
Because \(\mathcal S_0\) spans \(\widehat{\mathcal V}\), the
vectors \(B_s\phi\), \(s\in\mathcal U,\ \phi\in\mathcal S_0\),
span a subspace of dimension \(n\ge(1-4b^{-2})D\).
Select a basis \(\mathbf v\) of that subspace from these
vectors. Proposition~\ref{prop:short} and Hadamard's inequality
give
\[
 I\le 5D\log(1+b^2)B,\qquad
 \log\Delta_\sigma(\mathbf v)\le n h,\qquad
 h=H-m\mathcal R+L .
\]
Apply Proposition~\ref{prop:long}. Since \(D-n\le4b^{-2}D\)
and \(-h\le b^2x^{99}\) by \eqref{eq:coarse-main}, it gives
\begin{align*}
 D(H-3L)
 &\le nh+10D\log(1+b^2)B\\
 &=D(H-m\mathcal R+L)+(D-n)(-h)+10D\log(1+b^2)B\\
 &\le D(H-m\mathcal R+L)+4Dx^{99}+10D\log(1+b^2)B .
\end{align*}
After cancellation of \(DH\), we have
\[
 m\mathcal R\le 4L+4x^{99}+10\log(1+b^2)B\ll L .
\]
The constant is absolute for sufficiently large \(x\),
since \(a\ge xb\) and \(\log(1+b^2)\ll\log b\).
Divide by \(maS\) in \eqref{eq:two-weights} and use
\(f=\varphi(q)T\) and \(T/S\le\log b\). We obtain
\[
 R_{q,d}(x)\ll_q \frac{(\log b)^5}{m}
                 +|R_{q,d}(y)|+(\log b)|R_{1,d}(y)|.
\]
At the distant cutoff \(y\), the fixed-modulus prime number
theorem in \eqref{eq:pnt-tail} gives
\[
 |R_{q,d}(y)|+(\log b)|R_{1,d}(y)|\ll_{q,d,\delta}x^{-2}.
\]
Indeed \(\log y=x^{\delta/6}\), so a sufficiently high
logarithmic power in that estimate suffices. Also
\((\log b)^5/m\le x^{-1/2+5\delta}\).
Taking \(0<\delta<\min(1/100,\epsilon/5)\) proves
\eqref{eq:prime} for large \(x\).
The constant can be increased on a bounded interval
because \(R_{q,d}\) is continuous.
\end{proof}

The implication in \S\ref{sec:analytic} now proves
Theorem~\ref{thm:grh}, including the character of conductor
one.

\begin{thebibliography}{9}

\bibitem{alpoge}
L.~Alp\"oge,
Announcement of a proof of the Riemann hypothesis,
20 November 2026, \url{https://alpo.ge/remark.pdf}.

\bibitem{chojecki}
P.~Chojecki,
\emph{Covolumes and Dirichlet zeros},
preprint, 26 November 2026,
\url{https://www.ulam.ai/research/grh_complete.pdf}.

\bibitem{green}
J.~A. Green,
The characters of the finite general linear groups,
\emph{Trans. Amer. Math. Soc.} \textbf{80} (1955), 402--447.

\bibitem{mv}
H.~L. Montgomery and R.~C. Vaughan,
\emph{Multiplicative Number Theory I. Classical Theory},
Cambridge Studies in Advanced Mathematics, vol.~97,
Cambridge University Press, Cambridge, 2007.

\end{thebibliography}
\end{document}
